\begin{circuitikz}[
        american,
        voltage dir=RP,
        >={Latex[length=5mm, width=2mm]},
        %>=latex,
        alt={Schematic diagram of a simple DC electrical circuit connected to a PLC input point labeled Input Point. On the left, a DC battery labeled V provides voltage to the circuit. Connected in series from the positive terminal is a push button switch labeled PB. The circuit connects to a PLC input point shown inside a dashed red box. Inside the PLC input section, there is an internal current source representing a PNP sensor type input module. The lower return wire completes the circuit back to the negative terminal of the battery and is labeled Common.}
    ]

    \draw[
        BakerBlack
    ]
    (0,0)
    to[battery, color=BakerBlack, v={$V$}] ++(0,2)
    to[push button, color=BakerBlack, l={PB}, font=\small, -o] ++(3,0)
    coordinate(PLC_top)
    -- ++(1,0)
    to[isource, color=BakerBlack] ++(0,-2)
    to[short, color=BakerBlack, -o] ++(-1,0)
    coordinate(PLC_bottom)
    -- (0,0)
    node[
        pos=0.5,
        text=BakerBlack,
        font=\small,
        above
    ]{Common};

    \node[
        draw=BakerADARed,
        dashed,
        anchor=north west,
        minimum width=2cm,
        minimum height=3cm,
        label={above:{\small \textcolor{BakerBlack}{Input Point}}},
    ]
    (PLC)
    at ($(PLC_top)+(0,0.5)$)
    {};

    \node[
        text=BakerBlack,
        anchor=south
    ]
    at (PLC.south)
    {{\tiny PNP}};

\end{circuitikz}